\section{Pioneer SDK: полёт из Python}
\label{sec:python-sdk}

\begin{infobox}[Цели главы]
\begin{itemize}
    \item Запускать Python-программу в симуляторе, на настоящем дроне и с компьютера на симулятор.
    \item Подключаться к дрону, читать телеметрию и управлять подсветкой через \texttt{led\_control}.
    \item Выполнять взлёт и посадку безопасно: пауза после \texttt{arm()}, ожидание после \texttt{land()}.
    \item Летать по точкам и фигурам, дожидаясь каждой точки через \texttt{point\_reached()}.
\end{itemize}
\end{infobox}

\subsection{Где запускать Python}

Для управления Пионером из Python есть библиотека \texttt{pioneer\_sdk}: команды взлёта, посадки, полёта в точку, светодиоды, телеметрия, камера. У Python-программы, как и у Lua, два исполнителя.

\subsubsection*{В симуляторе}

Выберите в редакторе язык \ui{Python}. Библиотека уже «установлена»: \verb|from pioneer_sdk import Pioneer| работает, а дрон на сцене выполняет команды. Вывод \verb|print| --- в терминале, как и для Lua.

\begin{center}
\simshot[0.9\linewidth]{07-editor-python.png}
\captionof{figure}{Редактор в режиме Python. Переключение языка не теряет Lua-программу: у каждого языка своя вкладка.}
\end{center}

\subsubsection*{На настоящем дроне}

Python-программа работает на вашем компьютере и общается с дроном по Wi-Fi:

\begin{codebox}[Терминал]
\texttt{pip install pioneer-sdk}
\end{codebox}

Дрон создаёт сеть Wi-Fi (имя вида \texttt{Pioneer\_...}); подключитесь к ней --- и \verb|Pioneer()| найдёт дрон по адресу по умолчанию.

\subsubsection*{С компьютера --- на симулятор}
\label{sec:python-external}

Третий вариант: программа на вашем компьютере, а летает дрон в симуляторе. На сайте у дрона включается \ui{Разрешить внешние команды}, а в программе вместо адреса дрона указывается адрес сайта:

\begin{codebox}[Подключение к симулятору]
\begin{verbatim}
from pioneer_sdk import Pioneer
pioneer = Pioneer(ip='simulator.sakoryagin.ru',
                  mavlink_port=8001)
\end{verbatim}
\end{codebox}

Дальше программа ничем не отличается от программы для настоящего дрона. Так удобно отлаживать то, что в браузере недоступно: сторонние библиотеки, обработку видео с камеры.

\begin{infobox}[Lua ждёт событий, Python --- опрашивает]
В Lua следующую команду отдают из \verb|callback| по событию. В Python программа сама периодически спрашивает автопилот: \verb|while not pioneer.point_reached(): time.sleep(0.1)|. Смысл тот же --- не отдавать следующую команду, пока предыдущая не выполнена, --- меняется только форма.
\end{infobox}
